% CS224R Final Poster — Hybrid Advantage Shaping with Goal-Aware Attention TurnRD
% =============================================================================
% Compile with XeLaTeX (Overleaf: Menu > Settings > Compiler > XeLaTeX)
% Stanford LaTeX Poster Template (Gemini theme + Stanford color theme)
% 120cm x 72cm landscape, 3 columns, Palo Alto Green banner
% =============================================================================
\documentclass[final]{beamer}

% ============================================================================
% Packages
% ============================================================================
\usepackage[T1]{fontenc}
\usepackage{lmodern}
\usepackage[size=custom,width=120,height=72,scale=1.0]{beamerposter}
\usetheme{gemini}
\usecolortheme{stanford}
\usepackage{graphicx}
\usepackage{booktabs}
\usepackage{tikz}
\usepackage{pgfplots}
\usepackage{amsmath}
\usepackage{xcolor}
\pgfplotsset{compat=1.17}

% ============================================================================
% Lengths (3 columns)
% ============================================================================
\newlength{\sepwidth}
\newlength{\colwidth}
\setlength{\sepwidth}{0.025\paperwidth}
\setlength{\colwidth}{0.3\paperwidth}
\newcommand{\separatorcolumn}{\begin{column}{\sepwidth}\end{column}}

% ============================================================================
% Metadata (FILL IN AUTHOR NAMES + EMAILS)
% ============================================================================
\title{Hybrid Advantage Shaping with Goal-Aware Attention Credit Assignment\\for Multi-Turn LLM-Agent RL}

\author{Joseph Li \and Max Rodriguez \and Samantha Leventis}

\institute[shortinst]{~}

\footercontent{
  CS224R Final Project Poster Session \hfill
  \texttt{shoupei@stanford.edu, maxrod@stanford.edu, samanthaleventis@stanford.edu} \hfill
  Code: \texttt{github.com/JosephLi2023/CS224R}
}

% ============================================================================
% Body
% ============================================================================
\begin{document}

% Stanford logo top-right (template default)
\addtobeamertemplate{headline}{}{%
  \begin{tikzpicture}[remember picture,overlay]
    \node [anchor=north east, inner sep=3cm] at ([xshift=0.0cm,yshift=2.5cm]current page.north east)
      {\includegraphics[height=7.0cm]{stanford_logos/Block_S_2_color.png}};
  \end{tikzpicture}
}

\begin{frame}[t]
\begin{columns}[t]
\separatorcolumn

% =====================================================================
% COLUMN 1 — Motivation + Method + Architecture
% =====================================================================
\begin{column}{\colwidth}

  \begin{block}{Motivation}

    LLM agents acting over many turns in sparse-reward environments such as AlfWorld and
    WebShop receive only a single terminal reward, yet success depends on credit being
    assigned to the individual turns that actually mattered. As tool-use, web, and embodied
    agents tackle ever longer horizons, stable per-turn credit assignment has become the
    central bottleneck for sample-efficient learning.

    Existing approaches leave this gap open: trajectory-level GRPO (DeepSeekMath) copies one
    advantage onto every token and gives no per-turn signal; step-level GRPO (GiGPO, HoGPO)
    adds per-step structure but relies on fixed, heuristic credit; and return decomposition
    (RUDDER-style) learns per-step credit yet remains goal-agnostic. We therefore propose:

    \begin{itemize}
      \item \textbf{Hybrid Advantage Shaping (HAS)} --- a one-knob framework that blends a
            trajectory baseline with a pluggable per-turn decomposer.
      \item \textbf{TurnRD with goal-aware FiLM attention} --- the strongest decomposer:
            learnable attention over per-turn hiddens, modulated by a per-trajectory goal
            embedding.
    \end{itemize}

  \end{block}

  \begin{block}{Method: Hybrid Advantage Shaping}

    Per-turn advantage is a blend of the trajectory-level baseline and a per-turn signal:
    \begin{equation*}
      A_H^t \;=\; \alpha \cdot A_{\text{traj}} \;+\; (1-\alpha) \cdot A_{\text{turn}}^t
    \end{equation*}

    \begin{itemize}
      \item $\alpha = 1.0$: pure trajectory (flatGRPO, no credit);
            $\alpha = 0.5$: hybrid blend (empirically optimal)
      \item Per-turn decomposer is pluggable: Progressive,
            LLM-Judge, or TurnRD (learned attention)
    \end{itemize}

  \end{block}

  \begin{figure}
    \centering
    \IfFileExists{figures/architecture.png}{%
      \includegraphics[width=\linewidth]{figures/architecture.png}%
    }{%
      \IfFileExists{figures/F1_method_overview.pdf}{%
        \includegraphics[width=\linewidth]{figures/F1_method_overview.pdf}%
      }{%
        \fbox{\rule{0pt}{18cm}\rule{27cm}{0pt}}%
      }%
    }
  \end{figure}

  \begin{block}{TurnRD + FiLM (goal aware)}

    \begin{itemize}
      \item \textbf{Input/Encoder:} per-turn hiddens $h_t$ (policy's last hidden state) $\to$
            bidirectional transformer (2 layers, 4 heads, $H=128$)
      \item \textbf{FiLM (Feature-wise Linear Modulation) goal-conditioning:} per-trajectory goal\_emb $\to$ $\gamma, \beta$;
            $\text{FiLM}(h_t) = \gamma \odot h_t + \beta$ modulates $h_t$ before $\alpha$- and $V$-heads
      \item \textbf{$\alpha$-head:} softmax over turns $\to$ per-turn credit weight $\alpha_t$
      \item \textbf{$V$-head:} per-turn value $V_t$; $\sum_t \alpha_t V_t \approx R$ recovers reward
      \item \textbf{Training:} 4-term auxiliary loss (pred, rank, progress, value);
            re-fit per round on cumulative recency-decayed replay buffer
    \end{itemize}

  \end{block}

\end{column}

\separatorcolumn

% =====================================================================
% COLUMN 2 — AlfWorld Results (HEADLINE) + FiLM evidence
% =====================================================================
\begin{column}{\colwidth}

  \begin{block}{Experiment \& Results}

    \begin{itemize}
      \item \textbf{Headline results:}
      \begin{itemize}
        \item \textbf{AlfWorld: 80\%} success vs $\sim$70\% flatGRPO baseline (\textbf{+10pp})
        \item \textbf{WebShop: 91.3\%} success vs 80.5\% flatGRPO (\textbf{+10.8pp})
      \end{itemize}
      \item \textbf{Models:} Qwen2.5-1.5B-Instruct + LoRA $r=32$ (attn + MLP targets)
      \item \textbf{Envs:} AlfWorld valid\_seen, WebShop test split
      \item \textbf{Training:} 10 rounds $\times$ 80 episodes $\times$ $K=8$ rollouts/task
      \item \textbf{Eval:} greedy $K=1$, $n=100$ held-out episodes per round
      \item \textbf{Compute:} A100-80GB on Modal, $\sim$5--12h per run
    \end{itemize}

  \end{block}

  \begin{block}{AlfWorld Eval --- TurnRD+FiLM reaches 80\%}

    \begin{figure}
      \centering
      \includegraphics[width=0.85\linewidth]{figures/F4_alfworld_trajectory.pdf}
    \end{figure}

    \begin{itemize}
      \item \textbf{TurnRD+FiLM (80\%) > TurnRD no-FiLM (75\%) > flatGRPO ($\sim$70\%)}
            --- both contributions provide measurable lift
      \item Ours climbs smoothly and monotonically; flatGRPO regresses late
    \end{itemize}

  \end{block}

  \begin{alertblock}{Goal-Aware FiLM is Actively Learning}

    \begin{figure}
      \centering
      \includegraphics[width=\linewidth]{figures/F7_film_mechanism.pdf}
      \caption{Offline FiLM mechanism tests (3 panels): perturbation magnitude, goal-shuffle MSE, and $\gamma/\beta$ growth across rounds.}
    \end{figure}

    \begin{itemize}
      \item \textbf{Perturbation:} disable $\gamma/\beta$ $\to$ V-head output changes \textbf{39.7\%} --- FiLM is mechanically functional
      \item \textbf{Goal-shuffle:} shuffling goal\_emb $\to$ \textbf{+3.16\% MSE} ($57.6\%$ samples worse) --- FiLM is goal-aware
      \item \textbf{$\gamma/\beta$ growth:} $\gamma$ norm $0.34 \to 0.83$ monotone R0$\to$R12
    \end{itemize}

  \end{alertblock}

\end{column}

\separatorcolumn

% =====================================================================
% COLUMN 3 — WebShop + Conclusion + Limitations + Refs
% =====================================================================
\begin{column}{\colwidth}

  \begin{block}{WebShop --- HAS generalizes across envs}

    \begin{figure}
      \centering
      \includegraphics[width=\linewidth]{figures/F2_webshop_5method_bar.pdf}
    \end{figure}

    \begin{itemize}
      \item HAS @ $\alpha=0.5$ beats flatGRPO @ $\alpha=1.0$ by \textbf{+10.8pp} ---
            \textbf{the framework works}
      \item TurnRD attention beats Progressive (+5.3pp) and LLM-Judge (+9.3pp) ---
            \textbf{learned per-turn credit beats heuristics}
    \end{itemize}

  \end{block}

  \begin{block}{Conclusion}

    \begin{itemize}
      \item \textbf{HAS framework:} simple one-knob credit-assignment blend, generalizes across envs
      \item \textbf{TurnRD + FiLM:} strongest decomposer realization;
            \textbf{+5pp} additional lift from goal-aware FiLM on AlfWorld
      \item \textbf{80\% AlfWorld + 91.3\% WebShop} achieved with a single shared framework
      \item HAS at $\alpha = 1.0$ reduces to flatGRPO --- guaranteed never-worse-than-baseline
    \end{itemize}

  \end{block}

  \begin{block}{Limitations \& Future Work}

    \begin{itemize}
      \item \textbf{LoRA-only}: single base model (Qwen2.5-1.5B); full fine-tune untested
      \item \textbf{Two envs} (AlfWorld, WebShop); transfer behavior unknown
      \item \textbf{Weak goal-shuffle signal} ($+3.16\%$): goal-awareness is statistically real but modest in magnitude
      \item \textbf{Future:} adaptive per-round $\alpha$ schedule, larger base models,
            multi-task transfer, stronger goal-discrimination loss
    \end{itemize}

  \end{block}

  \begin{block}{References}

    \footnotesize
    \begin{thebibliography}{9}
      \setlength{\itemsep}{0.3em}
      \bibitem{feng2025gigpo} Yunzhen Feng et al. GiGPO: Step-level group relative policy optimization for long-horizon reasoning. arXiv:2505.10978, 2025.
      \bibitem{he2026hogpo} Shuo He et al. Hierarchy-of-groups policy optimization for long-horizon agentic tasks. In ICLR, 2026.
      \bibitem{shao2024deepseekmath} Zhihong Shao et al. DeepSeekMath: Pushing the limits of mathematical reasoning in open language models, 2024. Preprint.
      \bibitem{perez2018film} Ethan Perez et al. FiLM: Visual reasoning with a general conditioning layer. In AAAI, 2018.
    \end{thebibliography}

  \end{block}

\end{column}

\separatorcolumn
\end{columns}
\end{frame}

\end{document}
